\documentclass[a4paper,12pt,twoside]{report}

% 1. Voorkom 'Option clash' door opties vooraf te definiëren
\PassOptionsToPackage{backend=biber, style=numeric}{biblatex} 

% Pakketten
\usepackage[english]{babel} 
\usepackage[T1]{fontenc}
\usepackage[utf8]{inputenc} % Aanbevolen voor speciale tekens [cite: 9]
\usepackage[margin=2.5cm]{geometry}
\usepackage{amsmath}
\usepackage{graphicx}
\usepackage{booktabs}
\usepackage{hyperref}
\usepackage{enumitem}
\usepackage{microtype}
\usepackage{csquotes} % Verplicht voor biblatex i.c.m. babel 
\usepackage{../../school-macros}

% 2. BibLaTeX configuratie
\usepackage{biblatex} 
\addbibresource{referenties.bib} % De volledige naam incl. extensie [cite: 12, 71]

\begin{document}

% Titelinformatie
\kuleuventitle
{The relationship between religion and natural science: independent 
realms, irreconcilable conflict or complementary perspectives?}
{Religion}
{Louic Cooreman | R1002154}
{2025-2026}

\newpage
\tableofcontents
\newpage

\chapter{Introduction}
In this paper I will try to answer, in my own personal way, how I would define the relationship between religion and natural science. 
It is human for us to question the way these 2 topics are related just like it is human for us to agree or disagree on different topics, 
like for example politics. In this paper I will try to explain the conflict model of the relationship between natural science and religion, and it's limitations with an interesting example.
\newline
\newline
In the course book, \cite[Chapter 8, p. 344]{depoortere2023living}, it is mentioned that Ray Kurzweil predicts that : \textit{"artificial intelligence will surpass the brain capacity of humankind as a whole, and can be expected in 2045 and that eternal life, as an immortal human-machine hybrid"} is possible in the future. I will be discussing a similar topic to his prediction, if you have been keeping up with the news lately you may have seen articles or videos pop up about how scientists have placed a fruit fly's brain in a virtual body \cite{turkiyetoday2026fruitfly}.
\newline
\newline
% Een bron verschijnt alleen in de lijst als deze geciteerd wordt[cite: 102, 133].
% Gebruik \nocite{*} als je alle bronnen uit je .bib wilt tonen zonder citatie.
As an engineering student I find it important to keep up with the latest news and technology and when I read about this example I was fascinated but also scared, especially because of how close I believe Kurzweil's prediction is. 
Scientists already mapped the brain of a fruit fly in 2024 \cite{brainfacts2024connectome}, 
but giving the brain full access to a simulated body, and seeing how it perfectly replicates behaviours of an actual biological fruit fly, 
makes me personally raise some concerning questions like, "Are we all just biological machines?" or "Is the concept of having a soul now disproven?".
 These are obviously some harsh questions, which is one of the reasons why I wanted to dive into this topic more deeply. Another reason is that because I somewhat agree with the Conflict model of this relationship.

\chapter{The Conflict Model}
The fact that copying a brain of a biological creature, such as us, is now proven to be possible, 
makes it a perfect argument against the existence of a God (known as Materialistic atheism). 
\textit{"This is also known as an example of a conflict from the side of natural science".} \cite[Chapter 3, p. 76]{depoortere2023living}.
\newline
This argument for the incompatibility of natural science and religion is a conclusion of the conflict model of Jerry A. Coyne \cite[Chapter 3, p. 97]{depoortere2023living} 
which states that science and religion conflict with each other in a way that their method to understand reality is different. 
For science it's based on observations, proof and testing of theories (falsifiability), while religion focuses on belief without empirical proof. 
Another interesting fact to notice is that it's mentioned that \textit{"Coyne is a specialist in speciation and genetics (especially of \textbf{the fruit fly})".} 
He used this tiny creature as a perfect example because it's such a simple creature, also sometimes called a "model organism" \cite{doctrow2024completewiring}. 
And since this experiment showed that it is the brain itself that makes the decision and that reactions are purely mathematical it is safe to say that these reactions are explainable without the terms "soul" or "God".
\newpage
\subsection{The limitations of this relationship}
So why not just end the debate here? 
One could argue that, this experiment proves very little because it was done on a fruit fly, 
but I would argue that that is not a logical way of thinking because a fruit fly would still be part of God's creations. 
But maybe that doesn't even matter since the main goal of the company behind this simulation \cite{turkiyetoday2026fruitfly} is trying to simulate humans, 
fruit flies were just the start. But that is not the main issue with this confliction relationship. 
A bigger issue is that this experiment still did not prove that divinity like having a soul is false. 
Who is to say that it is not the soul who makes those mathematical decisions? 
So a major limitation of applying the conflict model to this experiment is that it mistakenly assumes it can disprove the existence of a "soul" or "God".
As discussed in \cite[Chapter 4, p. 192]{depoortere2023living} this is an example of \textit{"metaphysical ambiguity"}, 
which is perfectly illustrated by the \textit{"Parable of the invisible Gardener"} that \textit{Taede Smedes} mentioned \cite[Chapter 4, p. 186]{depoortere2023living}.
Shortly it concludes that people can look at the same reality with "different attitudes" and thus make different conclusions.
To further explore this, I wanted to hear the opinions of some of my peers so I asked them the following question : 
\newline
\newline
\textit{"If your brain could be copied and with it your entire personality would be simulated in a computer and u saw that it would make around 95 percent of the same conclusions or actions you would, would this make you believe More or Less in God?"}
\newline
\newline
Most of the people I asked just answered with a hard no, this is either because they are religious or aren't at all.
But other very interesting answers came up which I could agree with:
\newline
\newline
\textit{"I don't think this would be a reason to stop believing in God, also there's an emotional factor in us humans which I don't think we'll ever be able to simulate and even if we could, I don't think that simulating a person in a 'virtual world' would be proof of God not existing, I would however find it very cool if this is possible because I could kind of look at myself in the mirror and see what decisions I would take."}
\newline
\newline
When I heard what my peer said, I was surprised, mostly because of that "emotional factor" he talked about, not only because I agree, but also because this is exactly what Burms and De Dijn were talking about in \cite[Chapter 3, p. 84-90]{depoortere2023living} when they said that: \textit{"no description of biochemical processes, no matter how extensive, can ever capture what it actually means to be in love."}
\newpage
\chapter{Conclusion}
To conclude, do I still believe that religion and science conflict? Shortly, not necessarily. I think there is some truth to the limitations brought up in this paper. I could also argue about the complementary perspective because I believe that without religion, science wouldn't be as advanced as it is today, that is because I believe religion can be a "driving factor" for people to either try to prove or disprove religion using science, but obviously this isn't for everyone, I'm just saying that's a possibility.
But the main reason I believe they don't really conflict is because there's no concrete proof that God is real and so we should always keep in mind that if God is really that all-powerfull, well then he could definitely have made it so that it would always be up to debate.
\subsection*{Disclaimer:}
AI was used as a tool for translating, grammar error correction or shortly rewriting sentences, every thought in this paper is either mine, or referenced, so no additional text was generated. I also apologize for being slightly over the maximum of the 1000 words mark (excluding this disclaimer section).
%And who is to say that religion, more specifically the existence of God, was not one of the driving factors into starting this project of simulating brains in the first place? Would you spend your time trying to prove that we, as humans, are purely biomechanical creatures if you had solid evidence that God is real? 
%What about ethics, norms, morals? Are these not a part of religion? Who will decide in the future if it's morally okay for us to just make a copy of our mind and upload that into a simulation?
\newpage
\nocite{depoortere2023living}
\printbibliography[heading=bibnumbered] % Genereert de referentielijst [cite: 15, 348]

\end{document} % De absolute afsluiting van je bestand [cite: 17]